\documentclass[12pt,a4paper]{article}

\usepackage[utf8]{inputenc}
\usepackage[T1]{fontenc}
\usepackage[spanish]{babel}

\usepackage{geometry}
\usepackage{setspace}
\usepackage{array}
\usepackage{tabularx}
\usepackage{booktabs}
\usepackage{ragged2e}
\usepackage{float}
\usepackage{caption}
\usepackage{hyperref}
\usepackage[table]{xcolor}
\usepackage{tikz}
\usepackage{graphicx}
\usepackage{multirow}
\usepackage{amsmath}

\usetikzlibrary{babel,arrows.meta,positioning,shapes.geometric,calc,fit}

\geometry{margin=2.3cm}
\setstretch{1.15}
\captionsetup{font=small, labelfont=bf}

\hypersetup{
  colorlinks=true,
  linkcolor=blue!60!black,
  urlcolor=blue!60!black
}

\begin{document}

% ============================================================
%  ENCABEZADO
% ============================================================
\begin{center}
{\Large \textbf{Informe de Experimentos: Comparación de Algoritmos de\\
Machine Learning para Detección de Crisis Epilépticas en EEG Pediátrico}}\\[0.6cm]

\textbf{Proyecto:} Automatización de la detección de crisis en epilepsia focal pediátrica
mediante Inteligencia Artificial\\[0.3cm]

\textbf{Estudiante:} Erika Isabel Caita Avila\\
\textbf{Docente:} Alejandro Espejo\\
\textbf{Asignatura:} Investigación 2\\
\textbf{Programa:} Maestría en Inteligencia Artificial\\
\textbf{Universidad:} Universidad de La Salle\\
\textbf{Fecha:} Mayo de 2026\\[0.3cm]

\textbf{Repositorio del proyecto:}
\url{https://github.com/EkCaAv/Essentials-EEG}
\end{center}

\vspace{0.4cm}

% ============================================================
\section*{Resumen}
% ============================================================
\addcontentsline{toc}{section}{Resumen}

Este informe presenta el diseño, ejecución y análisis de un \textbf{experimento factorial
completo $4 \times 4$ con bloques aleatorios por sujeto} que compara cuatro algoritmos de
\textit{machine learning} clásico (Regresión Logística, Random Forest, SVM con kernel RBF
y Gradient Boosting) sobre cuatro conjuntos incrementales de características espectrales
y estadísticas extraídas de señales EEG. La cohorte está conformada por siete sujetos
pediátricos (6--10 años) del conjunto público \textbf{CHB-MIT Scalp EEG Database},
de los cuales se obtuvieron 52\,244 ventanas de 5 segundos con desbalance severo
(1,24\,\% ictales). La validación se realizó mediante \textit{GroupKFold} con $n=7$ folds
por sujeto, generando 112 entrenamientos en total. Los resultados muestran que dos
combinaciones superan el criterio de éxito mínimo del DOE (PR-AUC $>$ 0.15): Random
Forest con \texttt{bp\_plus\_rms\_kurt\_skew} (PR-AUC = 0.168 $\pm$ 0.166) y Random Forest
con \texttt{bp\_plus\_rms\_kurt} (PR-AUC = 0.153 $\pm$ 0.142). El test de Wilcoxon
\textit{signed-rank} pareado por fold no evidencia diferencias estadísticamente
significativas entre los modelos en ningún feature set ($p > 0.05$ en todos los pares),
hallazgo consistente con la baja potencia del test bajo $n=6$ folds válidos. El pipeline
es completamente reproducible y trazable, con código y artefactos disponibles públicamente
en el repositorio del proyecto.

\vspace{0.3cm}
\noindent\textbf{Palabras clave:} EEG pediátrico, epilepsia, detección de crisis,
Machine Learning, Random Forest, Gradient Boosting, GroupKFold, PR-AUC, CHB-MIT,
diseño factorial DOE.

\vspace{0.4cm}
\hrule
\vspace{0.3cm}

\tableofcontents

\vspace{0.4cm}
\hrule
\vspace{0.4cm}

% ============================================================
\section{Introducción}
% ============================================================

La epilepsia es uno de los trastornos neurológicos más frecuentes, con una prevalencia
cercana al 1\,\% de la población mundial \cite{who2019}. En población pediátrica, la
detección oportuna de crisis epilépticas a partir de registros de electroencefalografía
(EEG) resulta especialmente crítica, ya que las crisis no controladas generan daño
cognitivo acumulativo. Sin embargo, la revisión manual de registros de EEG de larga
duración es costosa, tardada y requiere especialización clínica avanzada.

El presente documento detalla el diseño, ejecución y análisis de un experimento factorial
comparativo que evalúa el desempeño de cuatro algoritmos de \textit{machine learning}
clásico —Regresión Logística, Random Forest, SVM con kernel RBF y Gradient Boosting—
sobre cuatro conjuntos de características espectrales y estadísticas extraídas de señales
EEG, en una subcohorte de siete sujetos pediátricos (6--10 años) del conjunto de datos
CHB-MIT. El experimento se enmarca en el segundo objetivo específico del proyecto de
investigación: analizar el desempeño de al menos dos técnicas de inteligencia artificial
aplicadas a la clasificación binaria ictal/interictal.

La justificación central de este enfoque comparativo reside en que no existe consenso
sobre qué representación de entrada y qué algoritmo produce el mejor balance entre
sensibilidad (detección de crisis) y especificidad (baja tasa de falsas alarmas) cuando el
desbalance de clases es severo. En los datos utilizados, apenas el 1,24\,\% de las
ventanas corresponde a actividad ictal (649 de 52\,244 ventanas), lo que invalida el uso
de exactitud (\textit{accuracy}) como criterio de evaluación y exige emplear el Área Bajo
la Curva Precisión-Sensibilidad (\textbf{PR-AUC}) como métrica primaria.

% ============================================================
\section{Diseño Experimental (DOE)}
% ============================================================

\subsection{Clasificación del diseño}

El experimento se clasifica como un \textbf{diseño factorial completo $4 \times 4$ con
bloques aleatorios por sujeto} \cite{montgomery2017}. La Tabla~\ref{tab:doe_clasif}
resume la estructura formal del diseño.

\begin{table}[H]
\centering
\caption{Clasificación del diseño experimental}
\label{tab:doe_clasif}
\renewcommand{\arraystretch}{1.25}
\begin{tabularx}{\textwidth}{>{\RaggedRight\arraybackslash}p{3.5cm}
                              >{\RaggedRight\arraybackslash}X}
\toprule
\textbf{Atributo} & \textbf{Descripción} \\
\midrule
Tipo de diseño     & Factorial completo $4 \times 4$ con bloques aleatorios \\
Factor A (tratamiento) & Algoritmo de ML: 4 niveles \\
Factor B (tratamiento) & Conjunto de características: 4 niveles \\
Factor de bloqueo  & Sujeto: 7 niveles (chb05, chb09, chb14, chb16, chb20, chb22, chb23) \\
Combinaciones      & $4 \times 4 = 16$ combinaciones de tratamiento \\
Entrenamientos totales & $16 \times 7 = 112$ \\
Estrategia de validación & \textit{GroupKFold} por sujeto ($n\_splits=7$) \\
Variable de respuesta primaria & PR-AUC (Average Precision) \\
Variable de respuesta secundaria & AUC-ROC, sensibilidad, especificidad \\
\bottomrule
\end{tabularx}
\end{table}

\subsection{Factores y niveles}

\paragraph{Factor A — Algoritmo de ML.}
Los cuatro algoritmos seleccionados cubren el espacio de hipótesis relevante para
clasificación tabular en contextos de desbalance severo:

\begin{table}[H]
\centering
\caption{Niveles del Factor A — Algoritmos de ML e hiperparámetros}
\label{tab:factor_a}
\renewcommand{\arraystretch}{1.3}
\begin{tabularx}{\textwidth}{>{\bfseries}p{3cm}
                              >{\RaggedRight\arraybackslash}p{5.2cm}
                              >{\RaggedRight\arraybackslash}X}
\toprule
\textbf{Algoritmo} & \textbf{Hiperparámetros clave} & \textbf{Justificación de elección} \\
\midrule
Regresión Logística &
  $C=1.0$, penalty=l2, class\_weight=balanced, solver=saga, max\_iter=20\,000 &
  Modelo lineal interpretable; actúa como \textit{baseline} regulado y permite
  estimar el límite inferior esperado de desempeño. \\
\midrule
Random Forest &
  n\_estimators=400, class\_weight=balanced\_subsample &
  Ensemble de árboles de decisión; estable ante alta dimensionalidad y desbalance
  por su submuestreo interno estratificado. \\
\midrule
SVM (RBF) &
  $C=10$, gamma=scale, class\_weight=balanced, probability=True, cache\_size=2\,000 &
  Clasificador de margen máximo; efectivo en espacios de alta dimensión normalizados.
  El kernel RBF captura relaciones no lineales entre características espectrales. \\
\midrule
Gradient Boosting &
  n\_estimators=300, lr=0.05, max\_depth=4, subsample=0.8 &
  Ensemble secuencial; manejo implícito del desbalance mediante pesos de muestra
  y mayor control del sobreajuste frente al RF. \\
\bottomrule
\end{tabularx}
\end{table}

\paragraph{Factor B — Conjunto de características.}
El diseño incremental permite evaluar el aporte marginal de cada tipo de característica:

\begin{table}[H]
\centering
\caption{Niveles del Factor B — Conjuntos de características}
\label{tab:factor_b}
\renewcommand{\arraystretch}{1.3}
\begin{tabularx}{\textwidth}{p{3.5cm} p{1.5cm} p{4cm} >{\RaggedRight\arraybackslash}X}
\toprule
\textbf{Nombre} & \textbf{N.\textsuperscript{o} feat.} & \textbf{Componentes} & \textbf{Justificación} \\
\midrule
\texttt{bp\_only} & 10 &
  Potencia relativa de bandas $\delta,\theta,\alpha,\beta,\gamma$ (media y desv. sobre canales) &
  Captura la distribución espectral baseline; evidencia de la literatura de que las
  bandas gamma y delta se ven afectadas durante las crisis. \\
\midrule
\texttt{bp\_plus\_rms} & 12 &
  Anterior + RMS por canal (media y desv.) &
  El RMS es una medida directa de la amplitud de la señal; las descargas ictales
  generan incrementos de energía detectables como picos de RMS. \\
\midrule
\texttt{bp\_plus\_rms\_kurt} & 14 &
  Anterior + Kurtosis por canal (media y desv.) &
  La curtosis mide impulsividad y colas pesadas; es sensible a los picos de alta
  amplitud y corta duración propios de las descargas epileptiformes. \\
\midrule
\texttt{bp\_plus\_rms\_kurt\_skew} & 16 &
  Anterior + Sesgo (\textit{skewness}) por canal (media y desv.) &
  El sesgo captura la asimetría de la distribución de amplitudes, que tiende a
  cambiar durante fases ictales con oscilaciones predominantemente positivas o negativas. \\
\bottomrule
\end{tabularx}
\end{table}

\subsection{Flujograma del experimento}

\begin{figure}[H]
\centering
\begin{tikzpicture}[
  scale=0.82, transform shape, >=Latex, font=\small,
  node distance=0.7cm and 1.2cm,
  block/.style={rectangle, rounded corners=3pt, draw=black, thick,
                fill=blue!8, text width=2.8cm, minimum height=0.9cm, align=center},
  decision/.style={diamond, draw=black, thick, fill=orange!15,
                   text width=2.2cm, align=center, inner sep=1pt},
  result/.style={rectangle, rounded corners=3pt, draw=black, thick,
                 fill=green!10, text width=3.0cm, minimum height=0.9cm, align=center},
  arrow/.style={->, very thick}
]
\node[block] (data)   {\textbf{Datos EEG}\\ CHB-MIT\\ 7 sujetos};
\node[block, right=of data]  (prep)   {\textbf{Preproceso}\\ Filtro 0.5--45 Hz\\ Resample 256 Hz};
\node[block, right=of prep]  (wind)   {\textbf{Ventaneo}\\ 5 s, 50\% overlap\\ 52 244 ventanas};
\node[block, right=of wind]  (feat)   {\textbf{Features}\\ BP, RMS,\\ Kurtosis, Skew};
\node[decision, below=1.2cm of feat] (fs) {Seleccion\\feature set\\(B1--B4)};
\node[block, left=1.5cm of fs]  (model)  {\textbf{Modelo ML}\\ (A1--A4)};
\node[block, left=of model]     (cv)     {\textbf{GroupKFold}\\ 7 folds\\ por sujeto};
\node[result, left=of cv]       (metric) {\textbf{Metricas}\\ PR-AUC\\ AUC-ROC};
\node[result, below=1.0cm of cv](out)    {\textbf{Resultados}\\ results.csv\\ cv\_fold\_metrics.csv};

\draw[arrow] (data)  -- (prep);
\draw[arrow] (prep)  -- (wind);
\draw[arrow] (wind)  -- (feat);
\draw[arrow] (feat)  -- (fs);
\draw[arrow] (fs)    -- (model)  node[midway,above,font=\scriptsize]{16 combos};
\draw[arrow] (model) -- (cv);
\draw[arrow] (cv)    -- (metric);
\draw[arrow] (metric)-- (out);
\draw[arrow] (out.south) .. controls +(0,-0.8) and +(0,-0.8) .. (fs.south)
  node[midway, below, font=\scriptsize]{siguiente combo};
\end{tikzpicture}
\caption{Flujograma del experimento. Cada una de las 16 combinaciones (Factor A $\times$
Factor B) se evalúa con 7 folds GroupKFold, produciendo 112 entrenamientos en total.}
\label{fig:flujo}
\end{figure}

% ============================================================
\section{Metodología}
% ============================================================

\subsection{Conjunto de datos: CHB-MIT Scalp EEG Database}

El conjunto de datos CHB-MIT \cite{chbmit} es una base pública del MIT y el Children's
Hospital Boston que contiene registros EEG de pacientes con epilepsia farmacorresistente.
Para este experimento se seleccionó la subcohorte pediátrica de siete sujetos entre 6 y 10
años (Tabla~\ref{tab:sujetos}), con el objetivo de mantener homogeneidad etaria y
consistencia con los objetivos del proyecto.

\begin{table}[H]
\centering
\caption{Sujetos incluidos en el experimento (recuentos efectivos del pipeline tras
preprocesamiento, ventaneo de 5 s con 50\,\% de solapamiento y submuestreo de
\texttt{max\_interictal\_per\_file=300})}
\label{tab:sujetos}
\renewcommand{\arraystretch}{1.2}
\begin{tabular}{llrrr}
\toprule
\textbf{ID sujeto} & \textbf{Archivos EDF} & \textbf{Ventanas ictales} &
\textbf{Ventanas interictales} \\
\midrule
chb05 & 39 & 0    & 11\,700 \\
chb09 & 19 & 117  & 5\,700 \\
chb14 & 26 & 82   & 7\,800 \\
chb16 & 19 & 52   & 5\,700 \\
chb20 & 29 & 131  & 8\,695 \\
chb22 & 31 & 87   & 9\,300 \\
chb23 & 9  & 180  & 2\,700 \\
\midrule
\textbf{Total} & \textbf{172} & \textbf{649} & \textbf{51\,595} \\
\bottomrule
\multicolumn{4}{l}{\small Total de ventanas: 52\,244 (1,24\,\% ictales). Conteos extraídos del} \\
\multicolumn{4}{l}{\small \texttt{run\_manifest.json} y \texttt{cv\_fold\_metrics.csv} de \texttt{doe\_v1\_all\_models}.}
\end{tabular}
\end{table}

La Figura~\ref{fig:dist_clases} evidencia el desbalance severo del dataset:
el 98,76\,\% de las ventanas son interictales y solo el 1,24\,\% son ictales.
Este desequilibrio motiva el uso de PR-AUC como métrica primaria y el parámetro
\texttt{class\_weight} en todos los modelos.

\begin{figure}[H]
\centering
\includegraphics[width=0.88\textwidth]{report_images/fig_distribucion_clases.pdf}
\caption{Distribución de clases en el dataset de 52\,244 ventanas. Izquierda: proporción
global ictal/interictal. Derecha: recuento real de ventanas por sujeto extraído del
\texttt{windows\_dataset.csv}; chb05 contribuye únicamente ventanas interictales, lo que
explica el fold con AUC indefinido (§3.4).}
\label{fig:dist_clases}
\end{figure}

\subsection{Preprocesamiento de señal}

Cada archivo EDF fue procesado con la biblioteca MNE \cite{mne2013} aplicando los
siguientes pasos en orden: (1) selección de canales EEG activos, (2) filtro paso-banda
Butterworth 0.5--45 Hz para eliminar deriva de línea base y ruido de alta frecuencia,
(3) remuestreo a 256 Hz para homogeneizar la tasa de muestreo entre sujetos y (4)
extracción de ventanas deslizantes de 5 segundos con 50\,\% de superposición
(\textit{overlap}).

\paragraph{Justificación de la longitud de ventana (5 s) y del solapamiento (50\,\%).}
La elección de \textbf{5 segundos} responde a un compromiso explícito entre tres factores:
(i) ventanas más cortas ($\leq$ 2 s) producen estimaciones espectrales con baja resolución
en frecuencia, especialmente en la banda $\delta$ (0.5--4 Hz); (ii) ventanas más largas
($\geq$ 10 s) atenúan la dinámica transitoria propia del inicio ictal y aumentan la
probabilidad de mezclar segmentos ictales e interictales en una misma ventana; (iii) la
literatura de detección de crisis sobre CHB-MIT \cite{chbmit} reporta resultados
consistentes para ventanas en el rango 1--10 s, siendo 5 s el valor más recurrente. El
\textbf{solapamiento del 50\,\%} se justifica como técnica de aumento efectivo de muestras
ictales (clase minoritaria) sin introducir sesgo significativo, ya que cada ventana
mantiene 2.5 s de información distinta respecto a la ventana adyacente.

\subsection{Extracción de características}

Para cada ventana temporal de $5 \times 256 = 1\,280$ muestras por canal, se extrajeron
cuatro tipos de características, cada una resumida como media y desviación estándar
sobre los canales activos del registro:

\begin{enumerate}
\item \textbf{Potencia relativa de banda (BP):} calculada mediante el método de Welch
  (densidad espectral de potencia), dividiendo la potencia de cada banda por la potencia
  total. Bandas: $\delta$ (0.5--4 Hz), $\theta$ (4--8 Hz), $\alpha$ (8--13 Hz),
  $\beta$ (13--30 Hz), $\gamma$ (30--45 Hz). Total: $5 \times 2 = 10$ características.

\item \textbf{RMS (Root Mean Square):} medida directa de la amplitud de la señal,
  sensible a cambios de energía durante las descargas ictales. Total: 2 características.

\item \textbf{Kurtosis:} momento estadístico de cuarto orden que cuantifica la
  impulsividad de la señal; valores elevados indican picos de corta duración y alta
  amplitud, característicos de las descargas epileptiformes. Total: 2 características.

\item \textbf{Sesgo (\textit{skewness}):} momento estadístico de tercer orden que mide
  la asimetría de la distribución de amplitudes. Total: 2 características.
\end{enumerate}

\subsection{Estrategia de validación cruzada}

Se empleó \textbf{GroupKFold con $n=7$ folds}, donde cada fold excluye completamente
un sujeto del conjunto de entrenamiento. Esta estrategia garantiza que no exista fuga
de información entre pacientes (\textit{data leakage}), lo cual es fundamental en
aplicaciones médicas donde el modelo debe generalizar a sujetos no vistos durante el
entrenamiento. La Figura~\ref{fig:flujo} ilustra este proceso.

\paragraph{Justificación de $n=7$ folds.}
La elección de \textbf{exactamente 7 folds} se deriva del número de sujetos de la cohorte
y obedece al esquema clínico estándar \textit{leave-one-subject-out} (LOSO): cada fold
deja un sujeto completo como conjunto de prueba y los seis restantes como entrenamiento.
Usar $n<7$ implicaría agrupar sujetos en el conjunto de prueba (perdiendo granularidad
para diagnosticar el desempeño por paciente), mientras que $n>7$ es imposible bajo la
restricción de \texttt{GroupKFold} (no puede haber más folds que grupos). Esta elección
maximiza la información estadística disponible (cada sujeto contribuye una vez como
prueba) y permite el análisis de variabilidad inter-sujeto que se presenta en §6.1.

\paragraph{Implicación del desbalance en la validación.}
El sujeto chb05 presenta cero ventanas ictales en el fold de prueba (fold=1/7), lo
que hace indefinida la métrica AUC en ese fold. Por ello, todos los experimentos
reportan \texttt{folds\_with\_nan\_auc = 1} y las medias se calculan sobre los
6 folds válidos restantes.

\paragraph{Regla de oro train/test.}
Todas las métricas reportadas en este informe
(Tablas~\ref{tab:resultados_completos} y~\ref{tab:wilcoxon_bprms}, y todas las
figuras de resultados) corresponden \textbf{exclusivamente} a las predicciones
sobre el fold de test de cada iteración \textit{GroupKFold}. Las métricas de
entrenamiento no se reportan: se utilizan únicamente como diagnóstico interno
de sobreajuste. Esta convención garantiza que ninguna cifra publicada esté
contaminada por información del fold de test durante el ajuste.

\subsection{Métricas de evaluación}

\begin{table}[H]
\centering
\caption{Métricas de evaluación y su justificación}
\label{tab:metricas}
\renewcommand{\arraystretch}{1.25}
\begin{tabularx}{\textwidth}{p{3.2cm} >{\RaggedRight\arraybackslash}X}
\toprule
\textbf{Métrica} & \textbf{Justificación} \\
\midrule
\textbf{PR-AUC} (primaria) &
  Robusta ante desbalance severo; captura la relación precisión--sensibilidad en la clase
  minoritaria (ictal). Es la métrica de decisión según el DOE. Criterio de éxito:
  PR-AUC $>$ 0.15 (baseline cap. 7: 0.112 $\pm$ 0.149). \\
AUC-ROC &
  Discriminación global; comparable con la literatura existente. No suficiente por
  sí sola ante desbalance. \\
Sensibilidad (Recall ictal) &
  Fracción de crisis detectadas; crítica en contexto clínico donde los falsos negativos
  (crisis no detectadas) tienen consecuencias graves. \\
Especificidad &
  Fracción de segmentos no-ictales clasificados correctamente; alta especificidad
  reduce las falsas alarmas. \\
\textit{Accuracy} &
  Referencia solamente; \textbf{no usarse como criterio principal} por el desbalance. \\
\bottomrule
\end{tabularx}
\end{table}

% ============================================================
\section{Desarrollo de los Experimentos}
% ============================================================

\subsection{Configuración técnica del pipeline}

El experimento se ejecutó con el script \texttt{pipeline/01\_chbmit\_experiments.py} sobre Python
3.13.13 con scikit-learn 1.8.0, MNE 1.12.1, NumPy 2.4.4 y SciPy 1.17.1.
El identificador del run fue \texttt{doe\_v1\_all\_models} y los artefactos se
almacenaron en \texttt{out\_thesis\_final/classical\_all\_models/}.

El pipeline ejecutó automáticamente las 16 combinaciones del diseño factorial, iterando
primero sobre los cuatro feature sets y luego sobre los cuatro modelos, produciendo
$16 \times 7 = 112$ entrenamientos. La Tabla~\ref{tab:pipeline_config} resume los
parámetros de ejecución. Los hiperparámetros de cada modelo (Tabla~\ref{tab:factor_a})
se fijaron \textit{a priori}; \textbf{no se realizó búsqueda anidada de hiperparámetros}
en este run (\texttt{tune\_hyperparams=false} en el manifiesto), por lo que el fold
de test nunca participa en su selección y se preserva la independencia entre el
ajuste del modelo y la evaluación.

\paragraph{Justificación de parámetros operativos.}
Dos parámetros del pipeline merecen mención explícita por su impacto sobre el análisis:

\begin{itemize}
\item \textbf{Umbral de decisión \texttt{threshold} = 0.5}: corresponde al punto de corte
  por defecto de la salida \texttt{predict\_proba}. Se mantiene fijo y \textbf{no se
  optimiza} porque las métricas primarias del DOE (PR-AUC y AUC-ROC) son métricas
  \textit{rank-based} que evalúan el ordenamiento de probabilidades sobre todos los
  umbrales posibles, no un umbral particular. Las métricas dependientes del umbral
  (sensibilidad y especificidad) se reportan únicamente como diagnóstico operativo.
\item \textbf{\texttt{max\_interictal\_per\_file} = 300}: limita el número de ventanas
  interictales muestreadas aleatoriamente por archivo EDF. Se justifica para
  (i) controlar el tiempo de cómputo (sin este límite el dataset alcanzaría
  $>$ 300\,000 ventanas), y (ii) atenuar parcialmente el desbalance extremo, pasando
  de un ratio $>$ 1:500 a $\sim$ 1:80, sin alterar la representatividad de los
  segmentos interictales (el muestreo es aleatorio uniforme sobre el archivo).
\end{itemize}

\begin{table}[H]
\centering
\caption{Parámetros de configuración del pipeline experimental}
\label{tab:pipeline_config}
\renewcommand{\arraystretch}{1.2}
\begin{tabular}{ll}
\toprule
\textbf{Parámetro} & \textbf{Valor} \\
\midrule
Run ID             & \texttt{doe\_v1\_all\_models} \\
Ventana temporal   & 5 segundos \\
Superposición      & 50\,\% (paso = 2.5 s) \\
Filtro paso-banda  & 0.5--45 Hz \\
Frecuencia de muestreo & 256 Hz \\
Estrategia CV      & GroupKFold($n=7$) \\
Semilla aleatoria  & 42 \\
Dataset resultante & 52\,244 ventanas, 21 columnas \\
Combinaciones      & 16 (4 modelos $\times$ 4 feature sets) \\
Entrenamientos totales & 112 \\
\bottomrule
\end{tabular}
\end{table}

\subsection{Disponibilidad del código y reproducibilidad}
\label{sec:reproducibilidad}

Para garantizar la \textbf{trazabilidad y reproducibilidad} del experimento, el código
fuente completo del pipeline, los scripts auxiliares, los manifiestos de ejecución y los
artefactos generados están disponibles de forma pública en el repositorio del proyecto
\cite{repo_essentials_eeg}:

\begin{center}
\url{https://github.com/EkCaAv/Essentials-EEG}
\end{center}

\paragraph{Estructura del repositorio.} El pipeline reproducible se organiza en cuatro
scripts numerados ubicados en el directorio \texttt{pipeline/}:

\begin{table}[H]
\centering
\caption{Scripts del pipeline reproducible}
\label{tab:pipeline_scripts}
\renewcommand{\arraystretch}{1.25}
\begin{tabularx}{\textwidth}{>{\ttfamily}p{5cm} >{\RaggedRight\arraybackslash}X}
\toprule
\textbf{Script} & \textbf{Función} \\
\midrule
01\_chbmit\_experiments.py & Extracción de \textit{features} EEG y entrenamiento
  de las 112 combinaciones (16 $\times$ 7 folds). \\
02\_consolidate\_results.py & Consolidación de la tabla comparativa y test de
  Wilcoxon \textit{signed-rank} pareado. \\
03\_generate\_report\_images.py & Generación automática de las 9 figuras del
  presente reporte. \\
04\_generate\_pdf\_report.py & Compilación del reporte PDF final. \\
run\_doe\_pipeline.ps1 & \textit{Runner} maestro: ejecuta los cuatro pasos en
  secuencia con un solo comando. \\
\bottomrule
\end{tabularx}
\end{table}

\paragraph{Comando de ejecución.} Una vez clonado el repositorio e instaladas las
dependencias (\texttt{pip install -r requirements.txt}), el experimento completo
se reproduce desde la raíz del proyecto con:

\begin{center}
\texttt{.\textbackslash pipeline\textbackslash run\_doe\_pipeline.ps1}
\end{center}

\noindent El tiempo de ejecución estimado es de \textbf{~45 minutos} en un equipo de
escritorio estándar. Tras finalizar, los artefactos siguientes quedan disponibles en
\texttt{out\_thesis\_final/}: \texttt{results.csv}, \texttt{cv\_fold\_metrics.csv},
\texttt{run\_manifest.json}, \texttt{comparison\_table.csv} y el PDF
\texttt{reporte\_final\_doe.pdf}.

\subsection{Descripción detallada de las 16 combinaciones}

Cada combinación se identifica por un par (feature set, modelo). A continuación se
describe el comportamiento observado en consola durante la ejecución, que evidencia
la variabilidad inter-sujeto característica de los datos de epilepsia pediátrica.

\paragraph{Combinaciones con \texttt{bp\_only} (Feature Set B1, 10 características).}
Este es el conjunto más simple, con únicamente la potencia relativa de las cinco bandas
espectrales. El modelo de Regresión Logística obtuvo AUC-ROC = 0.531 $\pm$ 0.203 y
PR-AUC = 0.054 $\pm$ 0.089, lo que coincide con un rendimiento cercano al azar,
coherente con la hipótesis del DOE de que \texttt{bp\_only} es insuficiente. Random
Forest y Gradient Boosting lograron AUC-ROC próximo a 0.62, pero con PR-AUC $<$ 0.07,
también por debajo del criterio de éxito. SVM obtuvo PR-AUC = 0.068. En todos los casos
se observó que fold=1/7 (sujeto chb05) produjo AUC = NaN por ausencia de positivos en
test.

\paragraph{Combinaciones con \texttt{bp\_plus\_rms} (Feature Set B2, 12 características).}
La incorporación del RMS produjo una mejora notable en Regresión Logística: AUC-ROC =
0.703 $\pm$ 0.199, PR-AUC = 0.112 $\pm$ 0.149. Este resultado replica exactamente la
línea base del capítulo 7 del proyecto (PR-AUC = 0.112 $\pm$ 0.149), lo que valida la
reproducibilidad del pipeline. Gradient Boosting alcanzó el mayor AUC-ROC de este bloque
(0.748), aunque con PR-AUC = 0.097, inferior a Regresión Logística. SVM presentó
variabilidad alta (AUC-ROC = 0.597 $\pm$ 0.215), indicando sensibilidad al perfil ictal
de cada sujeto.

\paragraph{Combinaciones con \texttt{bp\_plus\_rms\_kurt} (Feature Set B3, 14 características).}
La adición de kurtosis benefició principalmente a Random Forest (PR-AUC = 0.153), que
supera el criterio de éxito mínimo del DOE (PR-AUC $>$ 0.15). Gradient Boosting mejoró
su AUC-ROC (0.785) pero se quedó en PR-AUC = 0.149, levemente por debajo del umbral.
Regresión Logística no mejoró sustancialmente respecto al Feature Set B2, lo que sugiere
que la kurtosis aporta información no lineal que los modelos de árbol explotan mejor que
un modelo lineal.

\paragraph{Combinaciones con \texttt{bp\_plus\_rms\_kurt\_skew} (Feature Set B4, 16 características).}
El Feature Set completo benefició especialmente a Random Forest, que alcanzó
PR-AUC = 0.168 $\pm$ 0.166, el valor más alto del experimento. Gradient Boosting
obtuvo PR-AUC = 0.149 y AUC-ROC = 0.792, el AUC-ROC más alto del diseño. Regresión
Logística y SVM no mejoraron con la incorporación del sesgo, lo que indica que estos
modelos no aprovechan eficientemente la asimetría de la distribución de amplitudes.

% ============================================================
\section{Resultados}
% ============================================================

\subsection{Tabla completa de resultados}

La Tabla~\ref{tab:resultados_completos} presenta las métricas promediadas sobre los
7 folds GroupKFold para las 16 combinaciones del diseño factorial, ordenadas por
PR-AUC descendente.

\begin{table}[H]
\centering
\caption{Resultados completos — 16 combinaciones ordenadas por PR-AUC}
\label{tab:resultados_completos}
\renewcommand{\arraystretch}{1.25}
\small
\begin{tabular}{llrrrrr}
\toprule
\textbf{Feature Set} & \textbf{Modelo} &
\textbf{N feat.} & \textbf{PR-AUC} & \textbf{AUC-ROC} &
\textbf{Sensib.} & \textbf{Especif.} \\
\midrule
\rowcolor{green!12}
bp\_plus\_rms\_kurt\_skew & Random Forest     & 16 & 0.168 $\pm$ 0.166 & 0.732 $\pm$ 0.174 & 0.008 & 0.997 \\
\rowcolor{green!12}
bp\_plus\_rms\_kurt       & Random Forest     & 14 & 0.153 $\pm$ 0.142 & 0.760 $\pm$ 0.126 & 0.000 & 0.997 \\
bp\_plus\_rms\_kurt\_skew & Gradient Boosting & 16 & 0.149 $\pm$ 0.130 & 0.792 $\pm$ 0.157 & 0.019 & 0.991 \\
bp\_plus\_rms\_kurt       & Gradient Boosting & 14 & 0.149 $\pm$ 0.131 & 0.785 $\pm$ 0.149 & 0.016 & 0.990 \\
bp\_plus\_rms             & Reg. Logística    & 12 & 0.112 $\pm$ 0.149 & 0.703 $\pm$ 0.199 & 0.515 & 0.595 \\
bp\_plus\_rms\_kurt       & Reg. Logística    & 14 & 0.097 $\pm$ 0.131 & 0.692 $\pm$ 0.194 & 0.506 & 0.592 \\
bp\_plus\_rms             & Gradient Boosting & 12 & 0.097 $\pm$ 0.077 & 0.748 $\pm$ 0.157 & 0.012 & 0.988 \\
bp\_plus\_rms\_kurt\_skew & Reg. Logística    & 16 & 0.094 $\pm$ 0.128 & 0.688 $\pm$ 0.189 & 0.493 & 0.593 \\
bp\_plus\_rms             & Random Forest     & 12 & 0.083 $\pm$ 0.093 & 0.717 $\pm$ 0.143 & 0.000 & 0.997 \\
bp\_plus\_rms             & SVM (RBF)         & 12 & 0.083 $\pm$ 0.094 & 0.597 $\pm$ 0.215 & 0.038 & 0.985 \\
bp\_plus\_rms\_kurt       & SVM (RBF)         & 14 & 0.077 $\pm$ 0.085 & 0.638 $\pm$ 0.216 & 0.060 & 0.991 \\
bp\_plus\_rms\_kurt\_skew & SVM (RBF)         & 16 & 0.071 $\pm$ 0.077 & 0.643 $\pm$ 0.234 & 0.038 & 0.990 \\
bp\_only                  & SVM (RBF)         & 10 & 0.068 $\pm$ 0.052 & 0.634 $\pm$ 0.182 & 0.061 & 0.993 \\
bp\_only                  & Reg. Logística    & 10 & 0.054 $\pm$ 0.089 & 0.531 $\pm$ 0.203 & 0.417 & 0.556 \\
bp\_only                  & Gradient Boosting & 10 & 0.038 $\pm$ 0.027 & 0.617 $\pm$ 0.121 & 0.012 & 0.993 \\
bp\_only                  & Random Forest     & 10 & 0.032 $\pm$ 0.020 & 0.626 $\pm$ 0.117 & 0.000 & 0.998 \\
\bottomrule
\multicolumn{7}{l}{\small \textit{Filas sombreadas: las dos combinaciones que superan el criterio de éxito mínimo del DOE (PR-AUC $>$ 0.15)}} \\
\multicolumn{7}{l}{\small \textit{Métricas promediadas sobre 6 folds válidos (fold=1 con chb05 excluido por AUC indefinido).}}
\end{tabular}
\end{table}

\subsection{Mapas de calor — Vista factorial}

Las Figuras~\ref{fig:heatmap_prauc} y~\ref{fig:heatmap_aucroc} presentan los resultados
como matrices $4 \times 4$ que permiten identificar visualmente los efectos principales
y la interacción entre los dos factores experimentales.

\begin{figure}[H]
\centering
\includegraphics[width=0.82\textwidth]{report_images/fig_heatmap_prauc.pdf}
\caption{Mapa de calor de PR-AUC (media $\pm$ desv. estándar). Las celdas más oscuras
indican mejor discriminación para la clase ictal. Se observa que Random Forest y Gradient
Boosting con los feature sets B3 y B4 concentran los mejores valores.}
\label{fig:heatmap_prauc}
\end{figure}

\paragraph{Justificación del patrón observado en los heatmaps.}
La concentración de los mejores valores en la esquina inferior derecha (modelos de árbol
$\times$ feature sets ricos) tiene una causa identificable: las características de orden
estadístico superior (kurtosis y skewness) generan \textbf{relaciones no lineales y
fuertemente discriminativas} entre la clase ictal e interictal — los picos epileptiformes
producen colas pesadas en la distribución de amplitudes que se reflejan directamente en
estos momentos. Los modelos basados en árboles (Random Forest, Gradient Boosting) capturan
estas relaciones de manera natural mediante particiones recursivas del espacio de
características, mientras que (i) Regresión Logística asume una frontera lineal en el
espacio transformado, y (ii) SVM-RBF, aunque puede capturar no-linealidades, se ve
penalizado por la dimensionalidad efectiva y por el desbalance severo, que distorsiona
la estimación del margen óptimo.

\begin{figure}[H]
\centering
\includegraphics[width=0.82\textwidth]{report_images/fig_heatmap_aucroc.pdf}
\caption{Mapa de calor de AUC-ROC. Gradient Boosting con \texttt{bp\_plus\_rms\_kurt\_skew}
alcanza el mayor AUC-ROC del experimento (0.792).}
\label{fig:heatmap_aucroc}
\end{figure}

\subsection{Comparación por modelo y feature set}

\begin{figure}[H]
\centering
\includegraphics[width=\textwidth]{report_images/fig_barras_prauc.pdf}
\caption{PR-AUC por feature set y modelo (barras agrupadas, con desviación estándar).
La línea roja discontinua señala la línea base del DOE (0.112) y la línea naranja punteada
el criterio de éxito mínimo (0.15). Solo Random Forest con B3 (PR-AUC = 0.153) y B4
(PR-AUC = 0.168) supera ese umbral; Gradient Boosting con B3 y B4 se queda en 0.149,
levemente por debajo.}
\label{fig:barras_prauc}
\end{figure}

\begin{figure}[H]
\centering
\includegraphics[width=\textwidth]{report_images/fig_barras_aucroc.pdf}
\caption{AUC-ROC por feature set y modelo. Gradient Boosting y Random Forest dominan
en esta métrica, especialmente con los feature sets B3 y B4.}
\label{fig:barras_aucroc}
\end{figure}

\subsection{Balance sensibilidad--especificidad}

Un hallazgo relevante es la divergencia en el balance entre sensibilidad y especificidad
según el tipo de modelo (Figura~\ref{fig:sensib_especif}):

\begin{itemize}
\item \textbf{Regresión Logística} mantiene sensibilidades entre 0.42 y 0.52 con
  especificidades moderadas (0.55--0.60), lo que indica que intenta detectar crisis a
  costa de generar falsas alarmas.
\item \textbf{Random Forest, SVM y Gradient Boosting} presentan sensibilidades casi
  nulas (0.000--0.060) con especificidades muy altas (0.985--0.998), lo que indica que
  prácticamente clasifican todo como interictal. Aun así, generan mejores PR-AUC porque
  cuando predicen ictal, tienen alta precisión en los casos donde aciertan.
\end{itemize}

\paragraph{Justificación de la baja sensibilidad pese a \texttt{class\_weight=balanced}.}
La aparente paradoja —sensibilidad casi nula en RF, SVM y GB a pesar de tener configurado
el parámetro de balanceo de clases— se explica por el comportamiento del umbral fijo de
decisión (\texttt{threshold} = 0.5) sobre probabilidades calibradas para una distribución
extremadamente desbalanceada (1:80 después del muestreo, 1:500 sin él):

\begin{enumerate}
\item El parámetro \texttt{class\_weight=balanced} pondera las muestras durante el
  \textbf{ajuste} (función de pérdida), pero \textbf{no} reescala las probabilidades
  de salida para producir un umbral óptimo. La probabilidad media de la clase ictal
  predicha por estos modelos se mantiene baja ($<$ 0.10) porque la distribución
  empírica de la clase positiva es marginal.
\item Random Forest y Gradient Boosting promedian votos de árboles individuales: dado
  que la clase ictal es minoritaria, la mayoría de árboles vota interictal y el umbral
  por defecto raramente se supera.
\item SVM-RBF con probabilidades vía Platt scaling sufre el mismo efecto: las
  probabilidades calibradas tienden hacia la clase mayoritaria.
\item Regresión Logística es la excepción porque su salida probabilística no se promedia
  ni se calibra a posteriori, lo que produce predicciones más balanceadas a costa
  de mayor tasa de falsas alarmas (especificidad 0.55--0.60).
\end{enumerate}

\noindent Este hallazgo refuerza la decisión metodológica del DOE de utilizar PR-AUC
(rank-based, independiente del umbral) como métrica primaria, y motiva la recomendación
de explorar \textbf{calibración del umbral} y \textbf{sobremuestreo sintético (SMOTE)}
como trabajo futuro.

\begin{figure}[H]
\centering
\includegraphics[width=0.72\textwidth]{report_images/fig_sensibilidad_especificidad.pdf}
\caption{Diagrama de dispersión sensibilidad vs. especificidad. Cada punto representa
una combinación (modelo, feature set). Se observa la bifurcación entre modelos con alta
sensibilidad/baja especificidad (Regresión Logística) y alta especificidad/baja
sensibilidad (Random Forest, SVM, Gradient Boosting).}
\label{fig:sensib_especif}
\end{figure}

% ============================================================
\section{Análisis Estadístico}
% ============================================================

Antes de presentar los contrastes, conviene formalizar la distinción entre
\textbf{población objetivo}, \textbf{muestra disponible} y \textbf{representatividad},
ya que toda inferencia se interpreta a la luz de estos tres elementos.

\begin{table}[H]
\centering
\caption{Población, muestra y representatividad del experimento}
\label{tab:poblacion_muestra}
\renewcommand{\arraystretch}{1.25}
\begin{tabularx}{\textwidth}{>{\bfseries}p{3.2cm} >{\RaggedRight\arraybackslash}X}
\toprule
\textbf{Elemento} & \textbf{Descripción} \\
\midrule
Población objetivo &
  Niños con epilepsia focal de 6--10 años con registros EEG disponibles. \\
\midrule
Muestra disponible &
  7 sujetos del CHB-MIT en el rango de edad
  (chb05, chb09, chb14, chb16, chb20, chb22, chb23). \\
\midrule
Sesgo de muestra &
  La cohorte CHB-MIT presenta un 80\,\% de sujetos femeninos y una distribución
  bimodal de edades (predominantemente 2--4 años y 10--14 años); la subcohorte
  pediátrica seleccionada hereda este sesgo. \\
\midrule
Representatividad &
  \textbf{Limitada}: los hallazgos son válidos para esta cohorte. La generalización
  a la población objetivo requiere validación externa con sujetos no incluidos en
  CHB-MIT. \\
\bottomrule
\end{tabularx}
\end{table}

\subsection{Variabilidad inter-sujeto (distribución por fold)}

La Figura~\ref{fig:violin_folds} muestra la distribución de PR-AUC a nivel de fold para
cada feature set. La alta varianza observada en todos los modelos refleja la heterogeneidad
inter-sujeto: el desempeño varía de forma pronunciada dependiendo de qué sujeto se excluye
en cada fold. Este fenómeno es inherente a las señales EEG pediátricas, donde cada
paciente tiene un patrón ictal propio.

\begin{figure}[H]
\centering
\includegraphics[width=\textwidth]{report_images/fig_violin_folds.pdf}
\caption{Distribución de PR-AUC por fold (violinplot, $n=7$ folds) para cada feature set.
La línea horizontal roja indica la línea base del DOE (0.112). La alta dispersión confirma
que la variabilidad inter-sujeto es la principal fuente de incertidumbre experimental.}
\label{fig:violin_folds}
\end{figure}

\subsection{Test de Wilcoxon signed-rank}

Para evaluar si las diferencias en PR-AUC entre modelos son estadísticamente significativas,
se aplicó el \textbf{test de Wilcoxon signed-rank pareado por fold} \cite{wilcoxon1945}
a las 6 comparaciones posibles entre los 4 modelos, para cada uno de los 4 feature sets.
Este test no paramétrico es apropiado dado el tamaño reducido de la muestra
(efectivamente $n = 6$ folds válidos, ya que el fold con chb05 produce AUC indefinido
y se excluye del test) y la no normalidad esperada de las métricas.

\paragraph{Hipótesis pareadas.}
Para cada par de modelos $(i, j)$ con $i \neq j$ se contrastan, sobre los $n=6$ folds
válidos del feature set bajo análisis:
\begin{itemize}
\item $H_{0_{ij}}$: las distribuciones de PR-AUC por fold de los modelos $i$ y $j$
  son \textit{posiblemente} equivalentes (la mediana de las diferencias pareadas es
  cero).
\item $H_{1_{ij}}$: existe evidencia \textit{probable} de diferencia sistemática entre
  $i$ y $j$ a través de los sujetos evaluados.
\end{itemize}
La decisión se basa en el estadístico $W$ de Wilcoxon (suma mínima de rangos con
signo) y el $p$-valor exacto, con umbral $\alpha = 0.05$.

\paragraph{Advertencia de potencia y errores estadísticos.}
Con $n=6$ folds y $\alpha=0.05$, la potencia del test de Wilcoxon para efectos
medianos es aproximadamente 0.30--0.35 \cite{cohen1988}. Por tanto, los resultados
no significativos no demuestran equivalencia, sino \textbf{insuficiencia de evidencia}.
Bajo este tamaño muestral, el experimento acepta deliberadamente un riesgo de
\textbf{Error Tipo II} ($\beta$) elevado a cambio de mantener acotado el riesgo de
\textbf{Error Tipo I} ($\alpha=0.05$): un falso positivo —declarar superior a un
modelo cuando no lo es— se considera más costoso que un falso negativo en este
contexto exploratorio. Para complementar el $p$-valor, la Tabla~\ref{tab:wilcoxon_bprms}
reporta también el \textbf{tamaño del efecto rank-biserial} $|r| = |W^{+} - W^{-}| /
\binom{n+1}{2}$, donde $W^{+}$ y $W^{-}$ son las sumas de rangos con signo positivo
y negativo respectivamente. Por convención \cite{cohen1988}: $|r| < 0.1$ efecto
trivial, $0.1 \leq |r| < 0.3$ pequeño, $0.3 \leq |r| < 0.5$ mediano, $|r| \geq 0.5$
grande.

\begin{figure}[H]
\centering
\includegraphics[width=\textwidth]{report_images/fig_wilcoxon.pdf}
\caption{Matrices de p-valores del test de Wilcoxon signed-rank para los 4 feature sets,
representadas como \textbf{matriz triangular inferior} (la matriz es simétrica, por lo que
el triángulo superior es redundante; la diagonal en gris claro corresponde a
auto-comparaciones no aplicables). Celdas rojizas indican pares con menor p-valor; el
asterisco (*) marca $p < 0.05$. \textbf{Ningún par alcanza significancia estadística}
($p < 0.05$) en ningún feature set con $n=6$ folds válidos; los valores más bajos
($p \approx 0.0625$) corresponden a Gradient Boosting vs SVM y Random Forest vs SVM en
los feature sets B3 y B4, marginalmente cercanos al umbral.}
\label{fig:wilcoxon}
\end{figure}

Los resultados del test para el feature set de referencia \texttt{bp\_plus\_rms} son:

\begin{table}[H]
\centering
\caption{Resultados Wilcoxon signed-rank — feature set \texttt{bp\_plus\_rms}}
\label{tab:wilcoxon_bprms}
\renewcommand{\arraystretch}{1.25}
\begin{tabular}{lrrrp{4.4cm}}
\toprule
\textbf{Par de modelos} & \textbf{W} & \textbf{p-valor} & \textbf{$|r|$} & \textbf{Interpretación} \\
\midrule
Gradient Boosting vs Random Forest  & 6.0  & 0.4375 & 0.429 & n.s.; efecto mediano no resuelto por la potencia \\
Gradient Boosting vs Reg. Logística & 10.0 & 1.0000 & 0.048 & n.s.; efecto trivial \\
Gradient Boosting vs SVM            & 6.0  & 0.4375 & 0.429 & n.s.; efecto mediano no resuelto por la potencia \\
Random Forest vs Reg. Logística     & 10.0 & 1.0000 & 0.048 & n.s.; efecto trivial \\
Random Forest vs SVM                & 10.0 & 1.0000 & 0.048 & n.s.; efecto trivial \\
Reg. Logística vs SVM               & 8.0  & 0.6875 & 0.238 & n.s.; efecto pequeño \\
\bottomrule
\multicolumn{5}{l}{\small Test computado sobre los 6 folds válidos de \texttt{bp\_plus\_rms} (chb05 excluido por AUC indefinido).} \\
\multicolumn{5}{l}{\small $|r|$ = tamaño del efecto rank-biserial; convención Cohen: trivial $<$0.1, pequeño 0.1--0.3, mediano 0.3--0.5.} \\
\multicolumn{5}{l}{\small Ningún par alcanza $p < 0.05$; interpretación epistemológica según §6.2 (advertencia de potencia).}
\end{tabular}
\end{table}

% ============================================================
\section{Discusión y Justificación}
% ============================================================

\subsection{Efecto marginal de los feature sets}

La Figura~\ref{fig:efecto_marginal} permite evaluar la hipótesis $H_{0\_features}$
del DOE: \textit{``la adición de RMS, kurtosis o skewness no mejora PR-AUC respecto a
bandpower puro.''}

\begin{figure}[H]
\centering
\includegraphics[width=\textwidth]{report_images/fig_efecto_marginal.pdf}
\caption{Efecto marginal de la complejidad del feature set sobre PR-AUC (izquierda) y
AUC-ROC (derecha). La mejora es pronunciada para Random Forest y Gradient Boosting al
agregar kurtosis y skewness, mientras que SVM y Regresión Logística no se benefician
significativamente.}
\label{fig:efecto_marginal}
\end{figure}

La evidencia obtenida \textbf{rechaza parcialmente $H_{0\_features}$}: la adición de
kurtosis y skewness sí produce mejoras detectables en PR-AUC para Random Forest y
Gradient Boosting (de 0.083 a 0.168 y de 0.097 a 0.149, respectivamente), pero no
para Regresión Logística ni SVM. Esto sugiere una \textbf{interacción A$\times$B}:
el beneficio de las características estadísticas de orden superior depende del modelo
utilizado.

\subsection{Comparación con la línea base}

El criterio de éxito mínimo del DOE (PR-AUC $>$ 0.15) se cumple en dos de las 16
combinaciones (Tabla~\ref{tab:resultados_completos}): \texttt{bp\_plus\_rms\_kurt\_skew}
+ Random Forest (PR-AUC = 0.168) y \texttt{bp\_plus\_rms\_kurt} + Random Forest
(PR-AUC = 0.153). Las dos combinaciones de Gradient Boosting con los feature sets B3 y B4
obtuvieron PR-AUC = 0.149, levemente por debajo del umbral. La mejor combinación supera la
línea base del capítulo 7 (0.112) en $+0.056$ puntos absolutos, lo que representa una
mejora del 50\,\% relativa. Dado que la desviación estándar del baseline es 0.149, esta
mejora está dentro del rango de ruido experimental, lo que refuerza la advertencia de
potencia estadística del DOE.

\subsection{Limitaciones del experimento}

\begin{enumerate}
\item \textbf{Tamaño muestral reducido:} $n=7$ sujetos genera baja potencia estadística.
  Los resultados son válidos para esta cohorte pero no generalizables sin validación externa.
\item \textbf{Variabilidad inter-sujeto:} la desviación estándar de PR-AUC es del mismo
  orden que la media en todas las combinaciones, lo que indica que el efecto del sujeto
  domina sobre el efecto del tratamiento.
\item \textbf{Convergencia de Regresión Logística:} se observaron
  \textit{ConvergenceWarnings} en varios folds, incluso con \texttt{max\_iter=20\,000},
  lo que indica que el problema es inherentemente difícil para modelos lineales bajo
  desbalance severo. Los resultados de LogReg deben interpretarse con cautela.
\item \textbf{SVM con kernel RBF:} la sensibilidad casi nula (0.038) sugiere que el
  hiperplano de separación aprendido clasifica prácticamente todo como interictal,
  a pesar del parámetro \texttt{class\_weight=balanced}.
\end{enumerate}

% ============================================================
\section{Conclusiones}
% ============================================================

Los experimentos ejecutados bajo el diseño factorial completo $4 \times 4$ permitieron
extraer las siguientes conclusiones fundamentadas:

\begin{enumerate}
\item \textbf{El criterio de éxito mínimo del DOE se cumple.} Dos de las 16
  combinaciones superan PR-AUC $>$ 0.15: \texttt{bp\_plus\_rms\_kurt\_skew} + Random
  Forest (0.168 $\pm$ 0.166) y \texttt{bp\_plus\_rms\_kurt} + Random Forest
  (0.153 $\pm$ 0.142). Gradient Boosting con B3 y B4 alcanza 0.149, levemente
  bajo el umbral.

\item \textbf{Las características estadísticas de orden superior aportan valor.}
  La adición de kurtosis y skewness al feature set \texttt{bp\_plus\_rms} mejora el
  PR-AUC de Random Forest en $+102\,\%$ (de 0.083 a 0.168) y de Gradient Boosting
  en $+54\,\%$ (de 0.097 a 0.149). Respecto al baseline de bandpower puro
  (\texttt{bp\_only}), las mejoras son aún mayores (RF: +425\,\%; GB: +292\,\%).

\item \textbf{El efecto del modelo interactúa con el feature set.} Modelos de árbol
  (Random Forest, Gradient Boosting) explotan mejor las características no lineales
  (kurtosis, skewness) que modelos lineales o de margen (Regresión Logística, SVM).

\item \textbf{La validación GroupKFold por sujeto expone alta variabilidad inter-sujeto.}
  Las desviaciones estándar de PR-AUC son del mismo orden que las medias, lo que
  evidencia que la personalización de los modelos a cada paciente es una dirección
  de investigación prioritaria.

\item \textbf{Ningún par de modelos alcanza significancia estadística}
  ($p < 0.05$) en ningún feature set bajo el test de Wilcoxon \textit{signed-rank}
  con $n=6$ folds válidos. Los valores más cercanos al umbral son
  Gradient Boosting vs SVM y Random Forest vs SVM en los feature sets B3 y B4
  ($p \approx 0.0625$), marginalmente no significativos. Este resultado es
  atribuible principalmente a la baja potencia del test con tan pocos folds
  (potencia $\approx 0.35$), y no implica equivalencia real entre modelos.

\item \textbf{El pipeline es reproducible y trazable.} La ejecución genera
  \texttt{results.csv}, \texttt{cv\_fold\_metrics.csv} y \texttt{run\_manifest.json}
  con todos los parámetros necesarios para replicar el experimento exactamente.
\end{enumerate}

\vspace{0.5cm}

\noindent Como trabajo futuro se recomienda: (1) ampliar la cohorte a los 23 sujetos del
CHB-MIT para aumentar la potencia estadística, (2) explorar modelos de aprendizaje
profundo (CNN-1D, LSTM) sobre las mismas ventanas de señal cruda, y (3) evaluar
estrategias de sobremuestreo sintético (SMOTE) para atenuar el efecto del desbalance
severo sobre la sensibilidad.

% ============================================================
%  REFERENCIAS
% ============================================================
\begin{thebibliography}{9}

\bibitem{who2019}
  World Health Organization. (2019). \textit{Epilepsy: a public health imperative}.
  WHO Press.

\bibitem{montgomery2017}
  Montgomery, D.~C. (2017). \textit{Design and Analysis of Experiments} (9.ª ed.).
  Wiley.

\bibitem{chbmit}
  Shoeb, A.~H. (2009). \textit{Application of machine learning to epileptic seizure
  onset detection and treatment}. MIT PhD Thesis.
  Disponible en: \url{https://physionet.org/content/chbmit/1.0.0/}

\bibitem{mne2013}
  Gramfort, A. et al. (2013). MEG and EEG data analysis with MNE-Python.
  \textit{Frontiers in Neuroscience}, 7, 267.

\bibitem{wilcoxon1945}
  Wilcoxon, F. (1945). Individual comparisons by ranking methods.
  \textit{Biometrics Bulletin}, 1(6), 80--83.

\bibitem{cohen1988}
  Cohen, J. (1988). \textit{Statistical Power Analysis for the Behavioral Sciences}
  (2.ª ed.). Lawrence Erlbaum.

\bibitem{repo_essentials_eeg}
  Caita Avila, E.~I. (2026). \textit{Essentials-EEG: pipeline reproducible para
  detección de crisis epilépticas en EEG pediátrico}. Repositorio GitHub.
  Disponible en: \url{https://github.com/EkCaAv/Essentials-EEG}

\end{thebibliography}

\end{document}
